% ============================================================
%  THE ORTYX PROTOCOL
%  Part I — Immersion
%  Chapter 1: The Problem of Compression and the
%             Loss of Temporal Richness
% ============================================================

\chapter{The Problem of Compression and the Loss of Temporal Richness}
\label{ch:compression-loss}

\chapepigraph{We do not merely hear sound. We hear structure within
sound, and structure within structure. What compression removes is
not always what it appears to remove.}{Flyxion}

\noindent
Begin with something that is not in dispute.

Digital audio compression works by discarding information. The
question that practitioners have debated since the early days of
perceptual coding is which information can be safely discarded ---
meaning, which information the auditory system would not have
processed consciously in any case. The psychoacoustic models
underlying MP3, AAC, and their successors are built on careful
measurements of masking thresholds: frequencies that are rendered
inaudible by neighbouring louder frequencies, temporal components
too brief to register as distinct events, spatial information that
listeners cannot resolve. Discard what cannot be heard and the
bitrate drops dramatically with perceptually negligible loss.

This is a genuinely elegant solution to a real engineering problem,
and it has worked well enough that the compressed formats have become
the dominant medium through which most people now experience recorded
music. The engineering success is not in question. What is in question
--- and what the \Phoenix{} corpus makes its central concern --- is
whether the psychoacoustic models that define \enquote{what cannot be
heard} are drawing the right boundary.

The models were built around conscious perception. They identify what
listeners cannot report hearing, cannot discriminate in controlled
conditions, cannot use to identify compressed from uncompressed
material in blind tests. By those measures, they perform well. But
the question the Phoenix corpus raises is different: are there aspects
of temporal and harmonic structure that, while not consciously
discriminable in isolation, contribute to something broader than
discrimination --- to fatigue reduction, to immersive depth, to the
particular quality of attention that sustained musical engagement
requires? If the answer is yes, then perceptual coding models may
be correctly specifying one boundary while missing another.

This is not a fringe concern. Research in auditory scene analysis,
in the psychoacoustics of room acoustics, in the neuroscience of
auditory attention, and in the phenomenology of musical experience
has repeatedly suggested that the auditory system processes signal
structure at timescales and in spectral regions that do not register
as consciously attended content but that nonetheless influence the
character of the listening experience. The Phoenix corpus is
responding to a real and partially open question in perceptual
science. Understanding it requires taking that question seriously
before asking how well the corpus answers it.

\section{What Is Lost in Compression}
\label{sec:what-is-lost}

The standard account of lossy compression loss is spectral: certain
frequency components are removed or coarsely quantized. This account
is accurate as far as it goes, but it misses a dimension that becomes
important when thinking about music rather than speech intelligibility
or isolated tone discrimination. Music is not a collection of
frequency components that happen to co-occur in time. It is a
temporal structure within which frequency components stand in
relationships --- harmonic relationships, envelope relationships,
phase relationships, onset relationships --- that carry information
not reducible to their instantaneous spectral content.

Consider a plucked string. The attack transient contains a brief,
complex burst of energy spanning a wide frequency range. The
subsequent decay follows a pattern where different harmonics
decay at different rates, with the higher harmonics generally
decaying faster, producing the characteristic timbral evolution
that distinguishes a plucked string from a bowed one, from a
struck one, from a synthesized approximation. Perceptual coding
systems typically handle attack transients by switching to finer
time resolution during the onset and coarser time resolution during
the steady-state portion of the signal. They handle harmonic decay
by quantizing the spectral envelope at each analysis frame.

What these procedures preserve reasonably well: the pitch, the rough
timbral character, the dynamic envelope, the onset timing. What they
handle less cleanly: the precise relationships between how individual
harmonics evolve relative to one another over the course of the note.
The question is whether those relationships matter. In terms of
conscious discriminability under ABX conditions, often they do not.
In terms of what musicians and recording engineers report when
comparing high-quality transfers to compressed versions of the same
material, the judgments are more complex: something about the spatial
coherence, the sense of the instrument occupying a specific acoustic
space, the quality of decay, changes in ways that is difficult to
quantify but consistent across listeners.

\begin{epistemicstatus}{Phenomenological}
The claim that compression alters perceptual qualities beyond
conscious discriminability is grounded in listener reports and
practitioner judgments, not in controlled perceptual experiments
with the specific frameworks described in this monograph. The
observations are real and replicable at the level of informal
consensus; their mechanistic explanation remains open.
\end{epistemicstatus}

The Phoenix corpus takes this phenomenological observation as its
starting point and proposes a mechanistic account. The account
centres on what it calls \emph{spectral fracturing}: the disruption
not of individual frequency components but of the temporal
relationships between them --- the way harmonics evolve together,
the way onset transients establish initial harmonic structure, the
way decay patterns carry information about the physical source. The
claim is that these relationships are preserved in the original
signal, that they are selectively destroyed by lossy compression,
and that their destruction accounts for the perceptual qualities
that practitioners describe as \enquote{compression artifacts} even
in cases where standard psychoacoustic metrics suggest the loss
should be inaudible.

Whether this account is correct is a question that will require
careful unpacking. For now, it is sufficient to note that the
diagnostic gesture --- something about temporal relationship
structure, not just spectral content, is being lost --- is
consistent with existing literature and worth examining seriously.

\section{The Temporal Dimension of Auditory Structure}
\label{sec:temporal-dimension}

One of the conceptually interesting moves in the Phoenix corpus is
the systematic elevation of the temporal dimension of audio signals
over the spectral dimension. Conventional digital audio processing
thinks primarily in frequency space: the Fourier transform provides
the canonical representational basis, and operations on audio are
typically defined as manipulations of spectral content. The corpus
proposes an inversion: that temporal structure should be the primary
representational primitive, and that spectral information should be
derived from temporal patterns rather than the reverse.

This is not as radical a departure as it might appear. The auditory
system itself is better understood as a temporal processor than a
spectral analyser. The cochlea does perform a frequency decomposition,
but auditory nerve fibres encode information primarily through their
firing timing patterns, and the higher auditory pathway is exquisitely
sensitive to inter-aural timing differences, to the temporal fine
structure of sound within each frequency channel, and to the timing
relationships between onset events across frequency channels. The
auditory scene analysis literature, particularly the work following
from Albert Bregman's foundational studies, has established that
the perceptual grouping of sound into streams --- the ability to
hear distinct instruments in a complex mixture --- depends critically
on temporal coherence: components that start together, modulate
together, and share onset patterns tend to be grouped together.

The Phoenix corpus is drawing on this body of work, at least
implicitly, when it argues that temporal stability --- measured
through the jitter metrics of the \Marine{} --- is the primary
dimension along which salience should be assessed. A signal whose
timing is stable, whose peaks arrive at predictable intervals, whose
amplitude modulations follow a coherent pattern, is a signal that
the auditory system can group and track. A signal whose timing is
irregular, whose peaks are displaced from prediction, whose amplitude
modulations are incoherent, is a signal that resists grouping ---
that sounds, in experiential terms, degraded, flat, or diffuse.

The technical implementation of this insight is the \Marine{}, and
it is worth examining in some detail before any critical pressure is
applied.

\section{The Marine Salience Algorithm: First Encounter}
\label{sec:marine-first}

The \Marine{} is presented in the Phoenix corpus as a salience
detector that operates in the time domain, without recourse to the
Fast Fourier Transform or any global spectral decomposition. Its
operation is local and sequential: it examines a continuous
waveform peak by peak, computing at each peak a small set of
quantities that together constitute a salience estimate for that
moment in the signal.

The core quantities are straightforward. For each detected peak at
time $t_i$ with amplitude $a_i$, the algorithm computes the
instantaneous period $T_i = t_i - t_{i-1}$, the difference between
the current inter-peak interval and the previous one. It maintains
exponential moving averages of both period and amplitude:
\[
  \bar{T}_i = (1 - \alpha)\bar{T}_{i-1} + \alpha T_i, \qquad
  \bar{A}_i = (1 - \beta)\bar{A}_{i-1} + \beta a_i,
\]
where $\alpha$ and $\beta$ are smoothing parameters. From these
it derives period jitter $\Jp(i) = |T_i - \bar{T}_i|$ and
amplitude jitter $\Ja(i) = |a_i - \bar{A}_i|$, measuring how
much the current peak deviates from the recent local trend.

The salience estimate then combines local energy $E_i$, a harmonic
alignment score $H_i$ (measuring whether the current peak frequency
is consistent with a harmonic series), and a jitter-based coherence
term:
\[
  \SalFunc(i) = w_E E_i + w_H H_i
  + w_J\!\left(\frac{1}{1 + \Jp(i) + \Ja(i)}\right).
\]
The jitter term takes values near 1 when timing is highly stable
and approaches 0 as timing becomes increasingly irregular. The
weights $w_E$, $w_H$, $w_J$ are parameters of the system.

What is architecturally distinctive about this formulation is the
O(1) computational complexity per peak: the algorithm maintains
no global state beyond the running averages and requires no look-ahead.
This makes it suitable for real-time processing in resource-constrained
environments, a practical advantage that is separate from any
theoretical claims about what jitter measures. For embedded audio
systems, for low-latency processing pipelines, for hardware
implementations where FFT compute cost is prohibitive, a peak-level
salience estimator of this kind has clear engineering utility
regardless of how one evaluates the corpus's broader theoretical
architecture.

\begin{technote}[Technical Note: O(1) Salience Estimation]
The \Marine{} achieves constant-time-per-sample operation by
restricting its state to exponential moving averages over a sliding
window. The computational footprint per peak is:
two multiplications and additions (period and amplitude EMA
updates), two absolute value computations (jitter terms), one
reciprocal (coherence term), and three multiply-accumulates
(salience weighted sum). On modern hardware this is a handful of
clock cycles; on embedded processors it remains tractable without
dedicated DSP co-processors. The FFT-based alternative for a
comparable frequency resolution requires $O(N \log N)$ operations
per frame, where $N$ is the frame length, typically 512--8192
samples. For real-time applications with strict latency budgets,
the difference is significant.
\end{technote}

The conceptual inversion --- treating jitter as structural
information rather than noise to be suppressed --- is the move
that gives the algorithm its theoretical interest beyond the
engineering context. In conventional signal processing, jitter
is an imperfection to be minimized: clock jitter in digital systems
degrades performance, timing jitter in audio converters introduces
distortion. The Marine corpus reframes jitter as a measurement of
signal quality in a more fundamental sense: a signal with low
jitter at the peak level is a signal that is organized, structured,
and likely to carry the kind of temporal coherence that the
auditory system can group and track. A signal with high jitter
is one that has been disrupted --- whether by compression artifacts,
by room reflections, by dynamic range processing, or by the
underlying incoherence of noise.

This reframing is attractive because it aligns the algorithm's
primary variable with a genuine property of the signal rather than
a measurement artifact. Whether the alignment is tight enough to
support the full weight of the theoretical architecture built on
top of it is a question for later chapters.

\section{Compression as Temporal Disruption}
\label{sec:compression-temporal}

With the Marine salience framework in view, it becomes possible to
state the Phoenix corpus's central diagnostic claim in more precise
terms. Lossy compression, on this account, acts as a temporal
disruption mechanism. The window-based transform coding at the heart
of MP3 and AAC introduces discontinuities at block boundaries that
manifest as jitter at the peak level of the decoded waveform. The
psychoacoustic masking models that govern quantization decisions
operate on the spectral envelope within each block, without
reference to how peaks in the decoded waveform will relate to each
other across block boundaries. The result is that the spectral
content within each block may be well-preserved by psychoacoustic
standards while the temporal relationships between peaks across
blocks are systematically disrupted.

This is a specific, testable claim. It predicts that Marine-derived
jitter measurements should be systematically higher in
compressed-and-decoded audio than in the original waveform, even
in cases where the compressed version scores well on conventional
perceptual quality metrics. It predicts that this jitter increase
should be concentrated at block boundaries. It predicts that the
degree of jitter increase should correlate with the aggressiveness
of the compression --- with the bitrate, with the quantization
step size, with the degree of spectral coarsening applied.

These are predictions that could in principle be tested with
existing audio datasets and standard compression tools. The Phoenix
corpus does not, in the materials examined here, provide such a
test with the systematic controls that would establish the
prediction quantitatively. That gap will become important later.
For now, the claim is noted as a specific and in-principle
empirically accessible prediction, which already places it in
a better epistemic position than claims that resist operationalization
entirely.

\begin{epistemicstatus}{Operationally Grounded}
The prediction that lossy compression increases Marine-derived jitter
metrics at block boundaries is specific, measurable, and in principle
testable with existing tools. The prediction has not been formally
tested in the Phoenix corpus, but its structure is sound: it specifies
a measurable quantity, a direction of effect, a mechanistic rationale,
and a correlation with a known parameter (bitrate). This is the kind
of claim the \Ortyx{} audit process will classify as survivable pending
experimental confirmation.
\end{epistemicstatus}

\section{The Compression--Cognition Analogy}
\label{sec:compression-cognition}

The Phoenix corpus does not confine its framework to audio processing.
Having established the Marine salience metric as a tool for detecting
temporal coherence in audio signals, it extends the framework to a
much broader claim: that the relationship between salience-filtered
temporal structure and meaningful signal content is a general
principle of cognition, not a specific property of the auditory
modality.

The extension proceeds in several steps. First, the corpus notes
that the auditory system's preference for temporally coherent signals
is not unique to hearing: visual processing, proprioception, and
interoception all exhibit analogous stability preferences. Signals
that are temporally coherent --- that arrive in predictable patterns,
that modulate consistently, that fit within established temporal
frames --- are processed more efficiently, grouped more readily, and
experienced more richly than signals that are temporally chaotic.
Second, it observes that this stability preference is itself
structurally related to the efficiency of representation: a temporally
coherent signal is more compressible than an incoherent one, because
its future states are more predictable from its past states. Third,
it proposes that the relationship between coherence and
compressibility is not merely an engineering convenience but reflects
something about the structure of meaningful information: that meaning
tends to reside in structure that is stable enough to be tracked,
and that temporal stability is the primary marker of that structure.

This sequence of moves is interesting and partially supported by
existing cognitive science. The connection between predictive
processing accounts of perception and compression efficiency is
well-established in the theoretical neuroscience literature.
The claim that temporal coherence is a cross-modal salience marker
is consistent with work on sensory integration and attentional
selection. The stronger claim --- that temporal stability is the
\emph{primary} marker of meaningful structure across all modalities
and cognitive levels --- is where the corpus begins to outrun
its support.

But that observation belongs to a later chapter. The present task
is to follow the logic on its own terms, which means allowing the
extension to unfold before pressing on its joints.

\section{A Historical Note: When Constraint Forced Elegance}
\label{sec:historical-constraint}

Before examining the Phoenix corpus's specific proposals,
it is worth marking a historical observation that
will recur in Part~IV and Part~V. The tension between
rich signal content and representational economy is
not new. It was resolved, with striking elegance, by
hardware scarcity.

The MOS Technology SID chip --- the sound synthesis
engine of the Commodore 64 --- represents sound not
as sampled audio but as a procedural specification:
oscillator parameters, envelope sequences, filter
modulations, and timing registers that together constitute
a program for generating sound rather than a recording
of it. A minutes-long composition fits in a few kilobytes
not through lossy compression of a waveform but
through the replacement of the waveform by the
operator that generates it.

This is the distinction between extensional and
intensional representation --- a distinction that
Chapter~19 will formalize in detail. The SID developer
stored the structure, not the output. Constraint
forced this choice; but the choice turned out to be
correct in a deeper sense: the structure is more
reusable, more compact, and more honest about what
it is than the output would have been.

The Phoenix corpus's concern with what compression
destroys --- the temporal microstructure of acoustic
events --- can be read, at one level, as a lament
that modern audio codecs made the opposite choice
from the SID developer: they store a degraded waveform
rather than a generative operator. Whether this framing
survives the audit is a question for Parts~II and~III.
For now, it is sufficient to note that the intuition
the Phoenix corpus is reaching for --- that there is
something wrong with representing structure as output
--- has a genuine engineering precedent.

\section{A Framework That Wants to Unify}
\label{sec:unification-desire}

Every broad theoretical framework contains, at some level, a
unification ambition --- a drive toward discovering that apparently
disparate phenomena share a common underlying structure. This
ambition is epistemically double-edged. When the underlying structure
is real, unification produces genuine compression: it allows many
observations to follow from few principles, reduces the number of
independent explanatory commitments, and generates non-obvious
predictions. When the underlying structure is imposed rather than
discovered, unification produces the opposite: it inflates the
explanatory vocabulary to accommodate recalcitrant data, forces
disparate phenomena into a shared frame that distorts both, and
generates predictions that appear novel but are actually entailed
by the definitional choices embedded in the framework.

The Phoenix corpus has a strong unification ambition. It wants the
Marine salience metric to be simultaneously: a tool for audio quality
assessment, a model of auditory attention, a principle of cognitive
processing, a framework for AI architecture, and a philosophical
account of what it means for a signal to be meaningful. These
ambitions are stated progressively across the corpus, with each
layer appearing after the previous layer has been established enough
to lend support.

At the level of audio quality assessment, the ambition is plausible
and the engineering is concrete. At the level of auditory attention
modelling, the connection to existing literature is real if
incomplete. At the level of cognitive processing, the extension
requires assumptions that are not established by the audio work.
At the level of AI architecture, the framework is proposing a
research programme rather than reporting results. At the level of
philosophical account of meaning, the framework is making claims
that would require an entirely different order of argument to
support.

Whether this progressive extension constitutes a genuine scientific
programme --- one that generates testable predictions at each level
that could in principle disconfirm the framework --- or whether it
constitutes what will later be called Semantic Universality Collapse
($\FailGeo{3}$) --- the expansion of a framework's vocabulary until
it can assimilate any observation without constraint --- is the
central question of Parts~II and~III.

For now, the observation is recorded without verdict. The unification
ambition is noted as present, as partially grounded, and as in need
of careful level-by-level examination. The reader who finds the
framework exciting at this point --- who senses in it the possibility
of a genuinely unified account of signal, attention, cognition, and
meaning --- is responding to something real. The excitement is not
irrational. It is data about what the framework is doing well.
Whatever happens in the chapters that follow, that response deserves
to be taken seriously as a starting point.

\medskip

The next chapter turns to the most architecturally ambitious
component of the corpus: the \MEMeight{} wave-interference memory
system, which extends the salience framework from signal processing
into a full model of how memory and cognition might be grounded in
temporal dynamics rather than static representations. It is where
the framework becomes most original and, not coincidentally, where
the questions accumulate most rapidly.
